\documentclass[UTF8,fontset=none,11pt,a4paper]{ctexart}
\usepackage{ipara-handbook}
\usepackage{graphicx}
\handbooksetup{DS-08 图上的结构算法}{408 · 数据结构 DS-08}{Goal · Candidate · Commit · Invariant}
\begin{document}
\handbooktitle{图上的结构算法}{不要背七段代码：追踪候选域、提交时刻与仍然成立的事实}{408 · 数据结构 Topic DS-08}

\begin{corebox}{母问题：图上的“下一步”凭什么不会破坏全局目标？}
所有算法先压缩为同一条生成链：
\[
\boxed{\text{Goal}\to\text{Candidate Domain}\to\text{Local Action}
\to\text{Commit Point}\to\text{Invariant}\to\text{Global Result}}
\]
最容易混淆的不是循环语法，而是“候选对象是谁、何时可以永久确定”。
Prim 与 Kruskal 都求最小生成树，却分别在单棵树边界与全图森林之间选择；
Dijkstra、Bellman--Ford 与 Floyd 都做松弛，却分别固定顶点、路径边数上界与允许经过的中间点集合。
\end{corebox}

\begin{methodbox}{本册定位与 Owner 边界}
\begin{itemize}
  \item \kw{Owns：}Prim/Kruskal、Dijkstra/Bellman--Ford/Floyd、拓扑排序与 AOE 关键路径的目标、不变量、提交条件和失败证据。
  \item \kw{Uses：}DS07 提供图表示与遍历状态；Kruskal 直接复用 DS06 Union-Find；Prim/Dijkstra 复用 DS05 的“按优先级取候选”抽象契约。
  \item \kw{接口事实：}当前 Prim 用数组扫描，Dijkstra 用标准库最小优先队列，并未直接复用 DS05 的整数最大堆类；不能把概念依赖写成代码依赖。
  \item \kw{Stops：}堆、并查集和图存储内部机制不在本册重讲；工程路由协议只可调用最短路模型。
\end{itemize}
\end{methodbox}

\tableofcontents
\newpage

\section{先认候选域，再认算法名}

\begin{center}\small
\begin{tabularx}{\linewidth}{L{1.8cm}Y Y Y Y}
\toprule
算法 & 目标对象 & 候选域 & 提交时刻 & 失败证据 \\
\midrule
Prim & 一棵 MST & 当前树到树外的割边 & 树外顶点以最小跨边权加入 & 仍有顶点却无有限割边 \\
Kruskal & 一棵 MST & 全图尚未处理的边 & 最小边连接两个不同分量 & 处理完仍不足 $V-1$ 条边 \\
Dijkstra & 单源最短路 & 未过期的暂定距离 & 最小暂定距离出队 & 出现负权边 \\
Bellman--Ford & 单源最短路 & 每一轮的全部有向边 & 第 $i$ 轮保证至多 $i$ 条边的最短路 & $V-1$ 轮后仍可松弛 \\
Floyd & 全源最短路 & 所有点对 $(i,j)$ & 中间点集合由前 $k-1$ 个扩为前 $k$ 个 & 最终存在 $d[i][i]<0$ \\
拓扑排序 & 约束合法序列 & 当前入度为零的顶点 & 顶点出队并释放后继约束 & 输出顶点数小于 $V$ \\
关键路径 & AOE 最长工期与零时差活动 & 拓扑序上的事件和活动 & 正向、反向两遍完成后判时差 & 图有环或工期为负 \\
\bottomrule
\end{tabularx}
\end{center}

\begin{corebox}{统一记忆锚点}
不要记“某算法总是取最小”。先补完整句子：
\emph{从哪个候选域，按什么量取谁；取出后是暂时用于传播，还是从此永久确定？}
只要候选域或提交时刻答错，代码即使长得像模板，证明也已经失效。
\end{corebox}

\section{最小生成树：同一目标，两种生长几何}

\subsection{Prim：一棵树向外吞并}
设 $S$ 为已经进入树的顶点集合，\code{best[v]} 表示从 $S$ 跨到顶点 $v$ 的最轻边权。
每轮选择树外 \code{best} 最小的顶点加入 $S$，再用它的邻边缩小其他顶点的 \code{best}。

\begin{methodbox}{循环不变量}
已提交顶点始终连成一棵树；每个未提交顶点的 \code{best} 是当前割 $(S,V-S)$ 上到该点的最轻已知边。
割性质保证本轮最轻跨边是安全边。若树外仍有顶点而所有 \code{best} 都是无穷大，证据不是“程序没找到”，而是图不连通。

Prim 的初始化不是细节：任选起点令其 \code{best}=0，其余顶点为 $\infty$；每次提交顶点时记录使其进入树的 \code{parent} 边。若图有并列最轻跨边，不同 tie-break 可能得到不同但同权的 MST；若题目问唯一性，应检查是否存在某个割上的并列最轻边（或给出足以保证所有边权互异的条件）。数组扫描版每轮寻找最小未提交顶点，复杂度为 $O(V^2+E)$；堆版可降为 $O((V+E)\log V)$，但两者的割不变量相同。
\end{methodbox}

\lstinputlisting[language=C++,firstline=32,lastline=67,firstnumber=1,numbers=left,caption={Prim：候选域是当前树的边界}]{30_408/10_数据结构/08_图上的结构算法/code/ds08_graph_algorithms.hpp}

\subsection{Kruskal：多棵树逐步合并}
Kruskal 不维护“当前树”，而维护无环森林。边按权从小到大扫描；只有两端属于不同连通分量时，才把边提交并合并分量。

\begin{examplebox}{同一张图为什么会产生两种轨迹}
边为 $01(1),12(2),23(1),03(8)$。从顶点 $0$ 出发的 Prim 轨迹是
\[
\{0\}\to\{0,1\}\to\{0,1,2\}\to\{0,1,2,3\};
\]
Kruskal 则可先同时形成分量 $\{0,1\}$ 与 $\{2,3\}$，再用 $12(2)$ 合并。
两者总权都为 $4$，但状态几何完全不同：\kw{Prim 是单树扩张，Kruskal 是森林合并}。
\end{examplebox}

\lstinputlisting[language=C++,firstline=69,lastline=91,firstnumber=1,numbers=left,caption={Kruskal：Union-Find 维护森林分量}]{30_408/10_数据结构/08_图上的结构算法/code/ds08_graph_algorithms.hpp}

\begin{methodbox}{MST 的边界与并列权值}
Prim 与 Kruskal 的安全性来自割性质和环性质，而不是来自“每次看起来最小”。并列权值时具体边集可能不同，但只要图连通，所得树的总权值相同；若题目要求唯一的边集，必须再检查权值是否互异或给出 tie-break 规则。非连通图不能输出一棵覆盖全部顶点的 MST，只能得到最小生成森林，程序应把缺少 $V-1$ 条提交边作为可解释失败。
\end{methodbox}

\section{最短路：同一个松弛式，三种“已经解决”}

对边 $(u,v,w)$，三种算法共享局部动作
\[
d[v]\leftarrow\min\{d[v],d[u]+w\},
\]
但局部动作相同不代表外层循环可以互换。外层循环决定每一步之后究竟固定了什么。

\subsection{Dijkstra：固定当前最近的顶点}
非负权保证未来绕路只会继续增加长度。因此，当未过期的最小暂定距离 $(d,u)$ 出队时，不存在一条经未确定顶点绕回来使 $d[u]$ 更小的路径。

\begin{center}
\resizebox{0.94\linewidth}{!}{\providecolor{themebg}{HTML}{FAFAF7}
\providecolor{themefg}{HTML}{111111}
\providecolor{themegray}{HTML}{666666}
\providecolor{themecurve}{HTML}{1D63B8}
\providecolor{themealert}{HTML}{C53030}
\providecolor{themeamber}{HTML}{B86E00}
\begin{tikzpicture}[>=Stealth,font=\small,color=themefg,text=themefg]
\tikzset{
  v/.style={circle,draw=themefg,fill=themebg,minimum size=9mm,inner sep=0pt,line width=.9pt},
  fixed/.style={v,draw=themecurve,fill=themecurve!10!themebg},
  bad/.style={v,draw=themealert,fill=themealert!10!themebg},
  e/.style={->,line width=1.0pt,themegray},
  neg/.style={->,line width=1.2pt,themealert},
  card/.style={rounded corners=4pt,draw=themegray,fill=themefg!3!themebg,inner xsep=7pt,inner ysep=6pt,align=left}
}
\node[anchor=west,font=\bfseries\large] at (0,6.6) {Dijkstra：非负权让“当前最小暂定距离”可以永久提交};
\node[anchor=west,text=themegray] at (0,6.05)
  {左：非负权下，未提交区域再绕一圈只会增加代价；右：一条负边即可让已经提交的距离被后来路径改小。};

% left panel
\node[font=\bfseries,text=themecurve] at (3.7,5.1) {非负权：提交成立};
\node[fixed] (s1) at (1.4,3.8) {$s$};
\node[fixed] (a1) at (4.0,4.4) {$a$};
\node[v] (b1) at (4.0,2.8) {$b$};
\node[v] (c1) at (6.5,3.6) {$c$};
\draw[e] (s1) -- node[above,font=\scriptsize] {2} (a1);
\draw[e] (s1) -- node[below,font=\scriptsize] {5} (b1);
\draw[e] (a1) -- node[above,font=\scriptsize] {$\ge0$} (c1);
\draw[e] (b1) -- node[below,font=\scriptsize] {$\ge0$} (c1);
\node[card,anchor=north,text width=5.9cm] at (3.95,2.15)
  {$d[a]=2$ 是当前最小暂定值。任何经过尚未提交顶点的新路径，都只能在已有非负代价上继续增加，因此不可能再把 $a$ 降到 2 以下。};

\draw[themegray,line width=.5pt] (8.0,1.0)--(8.0,5.4);

% right panel
\node[font=\bfseries,text=themealert] at (12.1,5.1) {负权反例：提交失效};
\node[fixed] (s2) at (9.2,3.8) {$s$};
\node[bad] (a2) at (12.0,4.5) {$a$};
\node[v] (b2) at (12.0,2.7) {$b$};
\draw[e] (s2) -- node[above,font=\scriptsize] {2} (a2);
\draw[e] (s2) -- node[below,font=\scriptsize] {5} (b2);
\draw[neg] (b2) -- node[right,font=\scriptsize,text=themealert] {$-4$} (a2);
\node[anchor=west,font=\small] at (13.0,4.55) {先看到 $d[a]=2$};
\node[anchor=west,font=\small,text=themealert] at (13.0,3.85) {若此时永久提交 $a$};
\node[anchor=west,font=\small] at (13.0,2.85) {后来路径 $s\to b\to a$};
\node[anchor=west,font=\small,text=themealert] at (13.0,2.15) {代价 $5+(-4)=1<2$};
\node[card,anchor=north,text width=6.2cm] at (12.0,1.35)
  {失败点不是“松弛公式错了”，而是\textbf{提交证明}错了：未提交区域中的负边可以把未来绕路总价降下来。};

\node[anchor=west,font=\bfseries,text=themecurve] at (0,.2)
  {$\boxed{\text{非负边权}\Rightarrow\text{最小暂定值可提交};\quad\text{出现负边}\Rightarrow\text{立即停用该提交不变量}}$};
\end{tikzpicture}
}
\end{center}
右侧反例只改变一条边为负权：松弛动作本身仍合法，真正失效的是“当前最小暂定距离可以永久提交”的证明。

\begin{methodbox}{读代码时只盯三个状态}
\begin{enumerate}
  \item \code{distance[v]}：目前最好上界，不等于所有时刻都已最终确定；
  \item \code{pending}：可能含同一顶点的旧候选；
  \item \code{current\_distance != distance[current]}：丢弃过期候选，只有未过期最小项才承担提交语义。
\end{enumerate}
负权边会破坏“取出即确定”的证明，所以实现必须在运行前拒绝它，而不是算完再碰碰运气。

初始化时仅源点距离为 $0$，其余为 $\infty$，前驱为空；松弛前必须确认 $d[u]\ne\infty$，并用足够宽的整数承载加法。若采用邻接矩阵，扫描每个未确定顶点的全部候选邻居；若采用优先队列，允许同一顶点出现多个候选副本，但必须用“出队距离是否仍等于当前 distance”过滤旧项。路径恢复沿 predecessor 反向走到源点再逆序；不可达点没有合法路径，不能把 $\infty$ 当作可加的数。
\end{methodbox}

\lstinputlisting[language=C++,firstline=93,lastline=128,firstnumber=1,numbers=left,caption={Dijkstra：非负权、最小候选与过期项过滤}]{30_408/10_数据结构/08_图上的结构算法/code/ds08_graph_algorithms.hpp}

\subsection{Bellman--Ford：固定路径允许使用的边数}
第 $i$ 轮扫描全部边后，可以保证所有“至多含 $i$ 条边”的最短路已被传播。
没有可达负环时，简单最短路至多含 $V-1$ 条边；因此 $V-1$ 轮后若仍能松弛，就得到一个可达负环证据。

\begin{corebox}{最常漏掉的保护条件}
若 $d[u]=\infty$，边 $u\to v$ 不可参与加法，更不能据此更新 $v$。
“无穷大不参与算术”同时是语义规则和防溢出规则。
\end{corebox}

\lstinputlisting[language=C++,firstline=130,lastline=164,firstnumber=1,numbers=left,caption={Bellman--Ford：轮数就是路径边数上界}]{30_408/10_数据结构/08_图上的结构算法/code/ds08_graph_algorithms.hpp}

\subsection{Floyd：固定允许经过的中间点集合}
令 $D^{(k)}[i][j]$ 表示只允许前 $k$ 个顶点作为中间点时，$i$ 到 $j$ 的最短距离，则
\[
D^{(k)}[i][j]
=\min\left(D^{(k-1)}[i][j],
D^{(k-1)}[i][k]+D^{(k-1)}[k][j]\right).
\]
因此 \code{middle} 必须放在最外层：完成第 $k$ 层之后，整张矩阵才具有统一的“允许中间点集合”语义。

若需要恢复路径，通常维护 \code{next[i][j]} 或 \code{path[i][j]}：当经 $k$ 松弛成功时，令 $next[i][j]=next[i][k]$（或令前驱变为 $path[k][j]$），而不是把中间点 $k$ 直接当作下一跳。初始化时 $next[i][j]=j$ 仅对存在直接边的点对成立；无边点对保持空标记。最后若 $D[i][i]<0$，说明存在可达负权回路，相关最短路不是有限值；若只有负权边而无负环，Floyd 仍可正确处理。

\lstinputlisting[language=C++,firstline=166,lastline=198,firstnumber=1,numbers=left,caption={Floyd：最外层循环扩张允许中间点集合}]{30_408/10_数据结构/08_图上的结构算法/code/ds08_graph_algorithms.hpp}

\begin{methodbox}{三算法一眼分流}
\begin{itemize}
  \item \kw{Dijkstra：}固定顶点；需要非负权；适合单源。
  \item \kw{Bellman--Ford：}固定路径边数上界；允许负边；可检测源点可达负环。
  \item \kw{Floyd：}固定允许的中间点集合；处理所有点对；对角线负值暴露负环。
\end{itemize}
看到“松弛”还不能选算法，必须再问：题目要求固定哪一种状态？
\end{methodbox}

\section{DAG：从释放约束到识别零时差活动}

\subsection{拓扑排序：删除的不是点，而是前置约束}
入度为零表示一个顶点的全部前置约束已经释放。顶点出队后，对每条出边将后继入度减一；只有减到零的后继才进入候选队列。
若队列耗尽但仍有顶点未输出，这些顶点之间存在无法释放的循环依赖。

若某一步同时有两个或以上入度为零的候选，拓扑序不唯一；若每一步都恰有一个候选且最终输出 $V$ 个顶点，序列才唯一。这个判据直接来自 Kahn 的候选域，不需要枚举所有排列。DFS 版本可用三色状态：灰色节点再次遇到即为后向边和环，只有从后继全部完成后才把顶点压入结果栈，最后逆序得到拓扑序；不能把“DFS 访问顺序”直接当作拓扑序。

Kahn 算法中若把入度为零的候选放入栈、队列或优先队列，会得到不同的合法拓扑序：FIFO 倾向于较早发现的顶点，栈产生另一种深度偏序，最小堆可得到字典序最小序列。选择容器改变的是 tie-break，不改变“出队后删除其出边约束”的正确性。拓扑排序完成后应验证输出长度是否为 $V$；少于 $V$ 时剩余顶点的入度都无法降到零，环是可解释失败证据。

\subsection{DAG 表达式与共享子问题}
将表达式表示为 DAG 时，叶节点是操作数，内部节点是运算，边表示数据依赖。相同子表达式可共享一个节点，因而图的节点数不必等于表达式树展开后的节点数；求值必须按拓扑序，任何节点只有在所有前驱值已就绪后才能计算。若题目要求把 DAG 还原为中缀表达式，需要在分叉处明确括号和共享节点的重复打印策略；“拓扑序”只保证依赖合法，不保证算术运算符优先级。

\lstinputlisting[language=C++,firstline=200,lastline=230,firstnumber=1,numbers=left,caption={Kahn 拓扑排序：候选域是当前零入度点}]{30_408/10_数据结构/08_图上的结构算法/code/ds08_graph_algorithms.hpp}

\subsection{关键路径：正向求最早，反向求最迟}
AOE 网中，顶点是事件，边是有持续时间的活动。沿拓扑序正向做最长路松弛：
\[
ve[v]=\max_{u\to v}\{ve[u]+w(u,v)\}.
\]
项目工期是所有事件最早发生时间的最大值。再以工期初始化 $vl$，沿逆拓扑序反向收紧：
\[
vl[u]=\min_{u\to v}\{vl[v]-w(u,v)\}.
\]
对活动 $u\to v$，其最早开始与最迟开始分别为
\[
e=ve[u],\qquad l=vl[v]-w(u,v).
\]
只有 $e=l$，即总时差为零，活动才是关键活动。

\begin{center}
\resizebox{0.94\linewidth}{!}{\providecolor{themebg}{HTML}{FAFAF7}
\providecolor{themefg}{HTML}{111111}
\providecolor{themegray}{HTML}{666666}
\providecolor{themecurve}{HTML}{1D63B8}
\providecolor{themealert}{HTML}{C53030}
\providecolor{themeamber}{HTML}{B86E00}
\begin{tikzpicture}[>=Stealth,font=\small,color=themefg,text=themefg]
\tikzset{
  event/.style={circle,draw=themefg,fill=themebg,minimum size=11mm,inner sep=0pt,line width=.9pt},
  critical/.style={->,line width=1.6pt,themecurve},
  normal/.style={->,line width=.95pt,themegray},
  label/.style={font=\scriptsize,fill=themebg,inner sep=1.2pt},
  card/.style={rounded corners=4pt,draw=themegray,fill=themefg!3!themebg,inner xsep=7pt,inner ysep=6pt,align=left}
}
\node[anchor=west,font=\bfseries\large] at (0,7.0) {AOE 关键路径：事件算 $ve/vl$，活动算 $e/l$，零时差才构成关键链};
\node[anchor=west,text=themegray] at (0,6.45)
  {边上的前一个数是持续时间 $w$；括号内给出活动 $(e,l)$。蓝色实线表示 $l-e=0$ 的关键活动。};

\node[event] (v0) at (1.2,3.8) {$v_0$};
\node[event] (v1) at (4.3,5.2) {$v_1$};
\node[event] (v2) at (4.3,2.4) {$v_2$};
\node[event] (v3) at (8.0,3.3) {$v_3$};
\node[event] (v4) at (11.5,3.8) {$v_4$};

\draw[normal] (v0) -- node[label,above left] {$3\;(0,1)$} (v1);
\draw[critical] (v0) -- node[label,below left,text=themecurve] {$2\;(0,0)$} (v2);
\draw[normal] (v1) -- node[label,above] {$2\;(3,4)$} (v3);
\draw[critical] (v2) -- node[label,below,text=themecurve] {$4\;(2,2)$} (v3);
\draw[normal,bend left=15] (v1) to node[label,above] {$3\;(3,5)$} (v4);
\draw[critical] (v3) -- node[label,above,text=themecurve] {$2\;(6,6)$} (v4);

\node[anchor=south,font=\scriptsize,text=themegray,align=center] at (v0.north) {$ve=0$\\$vl=0$};
\node[anchor=south,font=\scriptsize,text=themegray,align=center] at (v1.north) {$ve=3$\\$vl=4$};
\node[anchor=north,font=\scriptsize,text=themegray,align=center] at (v2.south) {$ve=2$\\$vl=2$};
\node[anchor=south,font=\scriptsize,text=themegray,align=center] at (v3.north) {$ve=6$\\$vl=6$};
\node[anchor=south,font=\scriptsize,text=themegray,align=center] at (v4.north) {$ve=8$\\$vl=8$};

\node[card,anchor=west,text width=5.3cm] at (12.8,5.25)
  {\textbf{正向：事件最早时间}\\[-1mm]
   $ve[v]=\max(ve[u]+w)$。\\
   本例 $ve[v_3]=\max(3+2,2+4)=6$。};
\node[card,anchor=west,text width=5.3cm] at (12.8,3.15)
  {\textbf{反向：事件最迟时间}\\[-1mm]
   $vl[u]=\min(vl[v]-w)$。\\
   本例 $vl[v_1]=\min(6-2,8-3)=4$。};
\node[card,anchor=west,text width=5.3cm] at (12.8,1.0)
  {\textbf{活动判定}\\[-1mm]
   对 $u\to v$：$e=ve[u]$，$l=vl[v]-w$。只有 $e=l$ 才是关键活动。};

\node[anchor=west,font=\bfseries,text=themecurve] at (0,.55)
  {关键路径：$v_0\to v_2\to v_3\to v_4$，工期 $=8$。};
\node[anchor=west,font=\scriptsize,text=themegray] at (0,-.1)
  {图中纵向位置只为避让，不表示事件时间；真正时间只看 $ve/vl$ 数值与边权。};
\end{tikzpicture}
}
\end{center}
图中节点只承载事件的 $ve/vl$，边承载活动持续时间和 $(e,l)$；纵向排布只为避让，不表示事件时间。

若要估算缩短工期的代价，先找当前关键路径；非关键活动有正总时差，单独缩短它不会改变项目工期。只有缩短所有当前关键路径上的至少一条活动，工期才可能下降；存在多条互不共享的关键路径时，必须同时在每条路径上选活动。活动可缩短量受工期下限或资源上限约束，不能把“零时差”误当成“可无限压缩”。

对活动网络常同时计算四个量：事件最早发生时间 $ve$、事件最迟发生时间 $vl$，活动最早开始时间 $e$、活动最迟开始时间 $l$。关键活动的判据是 $l-e=0$；关键路径是从源事件到汇事件的一条零时差活动链，项目工期等于汇事件的 $ve$。若有多个零时差分支，关键路径可能不止一条，不能只输出任意一条就声称唯一。正向最长路必须在拓扑序上做，反向收紧必须在逆拓扑序上做；未先验证 DAG 时，两个循环都没有定义。

\begin{corebox}{为什么只算正向最长路还不够}
正向一遍只能回答“项目最早多久完成”；它不能判断某条活动能否延迟而不影响总工期。
关键活动判断必须依赖反向的最迟时间。代码若只返回 \code{earliest}，就没有完成关键路径算法。
\end{corebox}

\lstinputlisting[language=C++,firstline=232,lastline=281,firstnumber=1,numbers=left,caption={关键路径：最早、最迟与零时差边}]{30_408/10_数据结构/08_图上的结构算法/code/ds08_graph_algorithms.hpp}

\section{测试不是报喜，而是攻击不变量}

\lstinputlisting[language=C++,firstline=29,lastline=72,firstnumber=1,numbers=left,caption={同目标对照、失败边界与完整关键路径断言}]{30_408/10_数据结构/08_图上的结构算法/code/ds08_graph_algorithms_test.cpp}

测试按容易混淆的接口成对设计：Prim 与 Kruskal 在同图上应得到同一 MST 总权；Dijkstra 与 Bellman--Ford 在非负图上应得到相同单源距离；负边只拒绝 Dijkstra，负环才拒绝 Bellman--Ford；拓扑排序输出序列，关键路径还必须输出 $ve$、$vl$ 与关键边。

路径代价使用 \code{long long}，并用两条 $2\times10^9$ 的边验证结果为 $4\times10^9$，避免算法逻辑正确却被 32 位整数边界破坏。

\section{复杂度与适用边界}

\begin{center}\small
\begin{tabularx}{\linewidth}{L{2.3cm}C{2.4cm}Y Y}
\toprule
算法 & 本实现复杂度 & 正确性前提 & 显式失败 \\
\midrule
Prim & $O(V^2+E)$ & 无向连通带权图 & 非连通 \\
Kruskal & $O(E\log E)$ & 无向连通带权图 & 非连通 \\
Dijkstra & $O((V+E)\log V)$ & 有向图边权非负 & 负边 \\
Bellman--Ford & $O(VE)$ & 允许负边 & 源可达负环 \\
Floyd & $O(V^3)$ & 距离加法不溢出 & 任意负环 \\
拓扑排序 & $O(V+E)$ & 有向无环图 & 有环 \\
关键路径 & 当前实现 $O(VE)$ & AOE 网为 DAG，工期非负 & 有环、负工期 \\
\bottomrule
\end{tabularx}
\end{center}

\section{考场与代码复原协议}

遇到图算法题，按五步写在草稿上：
\begin{enumerate}
  \item \kw{Goal：}要树、单源距离、全源距离、合法顺序，还是总工期与关键活动？
  \item \kw{Candidate：}候选是跨割边、全局边、暂定距离、全部边、点对，还是零入度点？
  \item \kw{Commit：}本轮结束后，哪个对象从“候选”变为“永久成立”？
  \item \kw{Update：}更新距离、入度、分量、跨边权，还是最早/最迟时间？
  \item \kw{Failure：}负边、负环、非连通、有环、不可达与溢出分别怎样显式出现？
\end{enumerate}

\begin{corebox}{压缩复原}
只记五问：\kw{求什么？候选是谁？何时提交？提交后什么永真？什么证据让证明失效？}
当五问能在新图上独立回答，代码才真正被压缩成心智模型；只会背函数名和循环层数，还没有形成模型。
\end{corebox}

\begin{methodbox}{代码证据}
完整实现位于 \codepath{code/ds08_graph_algorithms.hpp}，断言测试位于
\codepath{code/ds08_graph_algorithms_test.cpp}。正文中的行号引用与当前 C++17 实现同步。
\end{methodbox}

\end{document}
